% !PRINTED 76
\weberpage{76}
%<<CONT
is at the same time contained in $\mathfrak{B}'$, then one can choose $a'$ so that $\chi_b(a')$ is not $= 1$. But then it follows from (9)
\weq{10}{(\chi_b, r) = 0;}
hence the theorem:

\gesperrt{1. If the period $\eta$ belongs not to $\mathfrak{A}$, but to a more comprehensive group $\mathfrak{A}'$, then every resolvent $(\chi_b, r)$ vanishes when $b$ is a number from $\mathfrak{B}$ that does not occur in $\mathfrak{B}'$.}

In order now to pursue further the conditions for this theorem, we must take into account the mode of formation of the characters.

We choose the notation as we introduced it at the close of \S.\,16. Then, if $\beta_{-1}, \beta_0, \beta_1, \ldots \beta_\mu$ are the indices of $b$ and $\nu_{-1}, \nu_0, \nu_1, \ldots \nu_\mu$ those of $n$,
\weq{11}{\chi_b(n) = \varepsilon_{-1}^{\beta_{-1}\,\nu_{-1}}\, \varepsilon_0^{\beta_0\nu_0}\, \varepsilon_1^{\beta_1\nu_1} \ldots \varepsilon_\mu^{\beta_\mu\nu_\mu}.}

If we denote by $r_0, r_1, \ldots r_\mu$ primitive roots of unity of the degrees $2^\lambda, q_1^{\varkappa_1}, q_2^{\varkappa_2}, \ldots q_\mu^{\varkappa_\mu}$, then we can represent every primitive $m$th root of unity in the form (Bd.\,I, \S.\,132)
\weq{12}{r = r_0 r_1 \ldots r_\mu,}
and if we determine $n_0, n_1, \ldots n_\mu$ from the congruences
\weq{13}{n \equiv n_0 \ (\mathrm{mod}\ 2^\lambda), \quad n \equiv n_1 \ (\mathrm{mod}\ q_1^{\varkappa_1}), \ldots n \equiv n_\mu \ (\mathrm{mod}\ q_\mu^{\varkappa_\mu}),}
then $(\chi_b, r)$ splits by (5) into the product of the following sums:
\weq{14}{\overset{n_0}{\Sigma}\, \varepsilon_{-1}^{\beta_{-1}\,\nu_{-1}}\, \varepsilon_0^{\beta_0\nu_0}\, r_0^{n_0}, \quad \overset{n_1}{\Sigma}\, \varepsilon_1^{\beta_1\nu_1}\, r_1^{n_1}, \ldots \overset{n_\mu}{\Sigma}\, \varepsilon_\mu^{\beta_\mu\nu_\mu}\, r_\mu^{n_\mu}.}

Now if $b$ is a number that occurs in $\mathfrak{B}$ but not in $\mathfrak{B}'$, then by theorem 1. one of these sums must vanish.

But for this, by \S.\,17, the necessary and sufficient condition is that one of the congruences
\weq{15}{\beta_0 \equiv 0\ (\mathrm{mod}\ 2), \quad \beta_1 \equiv 0\ (\mathrm{mod}\ q_1), \ldots \beta_\mu \equiv 0\ (\mathrm{mod}\ q_\mu)}
is satisfied, and indeed one whose modulus divides $m$ more than once.

This condition must first of all, when the hypothesis of theorem 1. holds, be fulfilled for every number $b$ that occurs in $\mathfrak{B}$ but not in $\mathfrak{B}'$. But if $b'$ is a number in $\mathfrak{B}'$, then for every integer exponent $x$ the product $b'^{x} b$ is in $\mathfrak{B}$, but
%CONT>>
